\documentclass[11pt,a4paper]{article}
\usepackage[utf8]{inputenc}
\usepackage[T1]{fontenc}
\usepackage[spanish,es-tabla]{babel}
\usepackage[margin=2.5cm]{geometry}
\usepackage{booktabs}
\usepackage{tikz}
\usetikzlibrary{arrows.meta,positioning}
\usepackage{listings}
\usepackage{xcolor}
\usepackage[hidelinks]{hyperref}

% permitir corte de linea en los guiones de \path (nombres de archivo largos)
\makeatletter\g@addto@macro\UrlBreaks{\do\-\do\\}\makeatother

\lstset{basicstyle=\ttfamily\small,breaklines=true,frame=single,
        backgroundcolor=\color{black!4},columns=fullflexible,
        literate={á}{{\'a}}1 {é}{{\'e}}1 {í}{{\'i}}1 {ó}{{\'o}}1 {ú}{{\'u}}1 {ñ}{{\~n}}1 {Á}{{\'A}}1 {Í}{{\'I}}1}

\title{OCR de 2014 en una sola PC del laboratorio:\\ directorios, comandos y qué esperar}
\author{}
\date{5 de noviembre de 2026}

\begin{document}
\maketitle

\begin{abstract}
Guía para ejecutar el OCR de los escaneos de 2014 en una única PC del laboratorio, sin
privilegios de administrador y sin internet. Todo lo necesario viaja en el USB: el
intérprete de Python, Tesseract con el idioma español y poppler son portátiles y se
ejecutan desde allí. El documento fija el directorio de trabajo, los comandos en orden, el
paralelismo para equipos con muchos núcleos, cómo comprobar el resultado y qué partes no
se han probado todavía.
\end{abstract}

\tableofcontents

%--------------------------------------------------------------------
\section{Punto de partida}

Los PDFs de 2014 están descargados completos y casi todos son fotografías de la
resolución sin capa de texto, así que \texttt{pdftotext} no devuelve nada y no pueden
convertirse a markdown hasta pasarles OCR.

\begin{center}
\begin{tabular}{lr}
\toprule
Concepto & Cantidad \\
\midrule
PDFs de 2014 en el manifiesto y en el USB & 11\,473 \\
Con texto de origen (ya convertidos a markdown) & 70 \\
Sin texto, que necesitan OCR & 11\,403 \\
Con OCR hecho durante las pruebas del 4--5 de nov.\ & 65 \\
\textbf{Pendientes aproximados} & \textbf{11\,338} \\
\bottomrule
\end{tabular}
\end{center}

El manifiesto (\path{manifiesto-pdfs.csv}) y el Excel limpio
(\path{todo-corte-suprema-completo-limpio.xlsx}) tienen 209\,180 filas cada uno y
coinciden con los PDFs de 2014 que hay en el USB. Se copiaron al USB y su copia es
idéntica a la original.

%--------------------------------------------------------------------
\section{Qué hay en el USB}

En esta guía \texttt{X:} es la letra que Windows asigne al USB en esa PC; se ve en el
Explorador de archivos.

\begin{lstlisting}
X:\
|-- ocr-offline\                   <-- directorio de trabajo
|   |-- parte-sola.cmd             todo 2014 en 3 tandas (el que se usa)
|   |-- ocr-lab.cmd                OCR directo sobre una carpeta, y -Convertir
|   |-- ocr-parte.ps1, ocr-lab.ps1 los scripts que lanzan los .cmd
|   |-- COMANDOS.txt               esta misma lista, en texto plano
|   |-- portable\                  Python, Tesseract (spa), poppler
|   |-- wheels\                    paquetes de Python, para reconstruir
|   `-- instaladores\              copias de los instaladores originales
|-- SCRAP\testeo\                  ocr.py, a_markdown.py y el resto del codigo
`-- corpus-corte-suprema\
    |-- pdfs\<especialidad>\2014\  los PDFs
    |-- markdown\                  markdown ya generado
    |-- manifiesto-pdfs.csv
    `-- todo-corte-suprema-completo-limpio.xlsx
\end{lstlisting}

\paragraph{No se instala nada.} Los \texttt{.cmd} llaman a PowerShell con
\texttt{-ExecutionPolicy Bypass} solo para ese proceso, y los scripts añaden Tesseract y
poppler al \texttt{PATH} únicamente de esa ventana. Al cerrarla no queda nada cambiado en
la PC.

%--------------------------------------------------------------------
\section{Procedimiento}

\subsection{Directorio donde entrar}

Siempre el mismo:
\begin{lstlisting}
X:\ocr-offline
\end{lstlisting}
Forma rápida: abrir esa carpeta en el Explorador, hacer clic en la barra de direcciones,
escribir \texttt{cmd} y pulsar Enter. Por consola: \texttt{X:} y luego
\texttt{cd \textbackslash ocr-offline}.

\subsection{Paso 1: prueba con 50 PDFs}

Se hace una vez, para comprobar que la PC deja ejecutar todo desde el USB y para medir
la velocidad real de esa máquina. Se hace directamente sobre el USB y sobre una sola
especialidad pequeña (civil, 820 PDFs), sin copiar nada al disco local:
\begin{lstlisting}
ocr-lab.cmd -Limite 50 -Especialidad civil
\end{lstlisting}
Debe terminar con \texttt{OCR hecho en 50 PDFs, 0 fallos}. El log imprime \texttt{s/doc},
los segundos efectivos por documento; con ese número se estima el total
(sección~\ref{sec:tiempos}). Para comparar configuraciones basta repetir el comando con
otros \texttt{-Workers} y \texttt{-Jobs} (sección~\ref{sec:paralelo}).

No conviene probar con \texttt{parte-sola.cmd -Limite 50}: el argumento se pasa a las tres
tandas y cada una copiaría antes todos sus PDFs al disco local.

\subsection{Paso 2: OCR de todo 2014 y conversión a markdown}

\begin{lstlisting}
parte-sola.cmd
\end{lstlisting}

El script trabaja en tres tandas, una detrás de otra. En cada una copia sus PDFs al disco
local, hace el OCR allí en paralelo, devuelve los PDFs al USB y borra la copia local. La
copia local evita escribir miles de veces en el USB durante el OCR. \textbf{Cuando
terminan las tres tandas convierte solo a markdown}, sin que haya que lanzar nada más (si
alguna tanda falla, el script se detiene y no convierte).

\begin{center}
\begin{tabular}{lp{8.2cm}rr}
\toprule
Tanda & Especialidades & PDFs & Tamaño \\
\midrule
A & penal & 3\,586 & 7,1 GB \\
B & contencioso administrativo laboral y contencioso administrativo previsional
    & 3\,895 & 12,1 GB \\
C & revisión de procedimiento coactivo, contencioso administrativo, civil, laboral,
    constitucional, familia civil, comercial, familia tutelar y familia penal
    & 3\,992 & 11,8 GB \\
\bottomrule
\end{tabular}
\end{center}

\paragraph{Espacio en disco.} Se necesita en el disco local al menos el tamaño de la
tanda más grande, \textbf{12,1 GB}, más un margen. El script comprueba el espacio libre
antes de copiar y se detiene con un mensaje si no alcanza. En ese caso, usar otro disco:
\begin{lstlisting}
parte-sola.cmd -Local D:\ocr
\end{lstlisting}
El USB tenía 21,8 GB libres al preparar el kit; el OCR reescribe los PDFs con un tamaño
casi igual (un PDF de 6 páginas creció un 1,4\%), así que no debería agotarlo.

\paragraph{Vigilar los primeros minutos.} El ciclo copiar, procesar y devolver no se ha
ejecutado nunca de principio a fin. Conviene mirar la tanda A hasta que aparezca el primer
\texttt{Devolviendo penal al USB}, antes de dejar la PC sola.

\paragraph{No desconectar el USB} mientras se ve \texttt{Devolviendo \ldots\ al USB}. Al
terminar todo, incluida la conversión, aparece \texttt{TERMINADO}.

\paragraph{Dónde queda cada cosa.} Todo dentro de \texttt{X:\textbackslash
corpus-corte-suprema\textbackslash}, en la carpeta de cada especialidad:
\begin{itemize}
  \item Los PDFs, con la capa de texto añadida, en el mismo sitio de origen:
        \path|pdfs\<especialidad>\2014\|.
  \item El markdown en \path|markdown\<especialidad>\2014\|, con la ruta que indica el
        manifiesto, igual que los demás años.
  \item Lo que no se pudo procesar: \path|pdfs\fallos-ocr.csv| (el OCR falló)
        y \path|markdown\necesitan-ocr-2014.csv| (siguen sin texto tras el
        OCR). Si el segundo está vacío o no tiene filas, quedó todo convertido.
\end{itemize}

\subsection{Si se corta}

Un apagón o una ventana cerrada no estropean los PDFs: la escritura es atómica (archivo
temporal y renombrado). Basta lanzar \emph{exactamente} el mismo comando. Cada PDF que ya
tiene texto se salta.

\subsection{Conversión a markdown por separado}

\texttt{parte-sola.cmd} ya la hace al final. Solo hace falta lanzarla a mano si falló o
si se quiere repetir (por ejemplo, tras corregir PDFs en otra PC):
\begin{lstlisting}
ocr-lab.cmd -Convertir
\end{lstlisting}
Esto ejecuta \texttt{a\_markdown.py} para 2014 con el manifiesto y el Excel del USB y deja
el resultado en \texttt{X:\textbackslash corpus-corte-suprema\textbackslash markdown\textbackslash}.
El script detecta por fecha qué PDFs son más nuevos que su markdown y solo rehace esos.

%--------------------------------------------------------------------
\section{Paralelismo}\label{sec:paralelo}

Con muchos núcleos rinde más procesar varios documentos a la vez que repartir las
páginas de uno solo, porque la mayoría de los PDFs tiene pocas páginas (unas 7 a 10 de
media en las muestras medidas). El OCR se parametriza con dos números:

\begin{center}
\begin{tabular}{lp{9cm}}
\toprule
Parámetro & Significado \\
\midrule
\texttt{-Workers} & Documentos procesados a la vez. Por defecto: núcleos/2, máximo 16. \\
\texttt{-Jobs} & Hilos por documento. Por defecto: núcleos/\texttt{Workers}. \\
\bottomrule
\end{tabular}
\end{center}

En una PC de 24 núcleos el valor automático es de 12 documentos a la vez con 2 hilos cada
uno. Además, Tesseract se limita a un hilo por proceso (\texttt{OMP\_THREAD\_LIMIT=1}),
para que los procesos no compitan entre sí, y la revisión previa de qué PDFs ya tienen
texto también se hace en paralelo.

Cada documento en curso usa del orden de 0,5 a 1 GB de memoria. Si la PC se queda sin
RAM, bajar \texttt{-Workers}. Los parámetros se añaden al final de cualquier comando:
\begin{lstlisting}
parte-sola.cmd -Workers 12 -Jobs 2
parte-sola.cmd -Workers 6
parte-sola.cmd -Workers 1 -Jobs 24
\end{lstlisting}
La última línea es el modo anterior: un documento cada vez con todos los hilos. En la
prueba de 50 conviene comparar el \texttt{s/doc} de dos configuraciones y quedarse con la
mejor.

%--------------------------------------------------------------------
\section{Cuánto tardará}\label{sec:tiempos}

\paragraph{Lo medido.} En la PC de casa (4 núcleos), con un documento a la vez y 4 hilos,
50 PDFs de civil (520 páginas) tardaron unos 21 s por documento, es decir, unos 2,0 s por
página. Con esa velocidad, los 11\,338 pendientes serían unas 66 horas.

\paragraph{Lo que no se ha medido.} En la máquina de 24 núcleos no se ha ejecutado nada. Si
el rendimiento escalara linealmente con los núcleos, serían del orden de 11 horas, pero eso
es una cota optimista: el disco, la memoria y el propio Tesseract no escalan perfectamente.
Un intento de comparar en la PC de casa 4 documentos en paralelo contra uno solo salió peor
(4,0 contra 2,0 s por página), pero se hizo con documentos distintos, sobre un equipo
pequeño y con dos hilos reales, así que no permite concluir nada sobre el laboratorio. La
cifra útil saldrá de la prueba con 50 PDFs en la propia máquina.

\paragraph{Revisión previa.} Antes de empezar cada tanda, \texttt{ocr.py} recorre sus PDFs
para ver cuáles ya tienen texto. En la PC de casa, con la versión secuencial, recorrer
los 11\,473 de 2014 tardó unos 30 minutos; ahora se hace en paralelo y debería ser bastante
menos, pero no se ha medido sobre el conjunto completo.

%--------------------------------------------------------------------
\section{Comprobar el resultado}

\begin{itemize}
  \item El log queda en \texttt{X:\textbackslash corpus-corte-suprema\textbackslash
        ocr-lab-2014-NOMBREPC.log} y cada tanda debe terminar con \texttt{0 fallos}.
  \item Si algún PDF falla, aparece en
        \texttt{X:\textbackslash corpus-corte-suprema\textbackslash pdfs\textbackslash
        fallos-ocr.csv} con el motivo.
  \item \texttt{ocr.py} da por bueno un PDF solo si el resultado devuelve al menos 200
        caracteres con \texttt{pdftotext}; si no, conserva el original.
  \item Para saber qué queda pendiente basta lanzar \texttt{ocr-lab.cmd}: recorre los PDFs
        del USB e imprime \texttt{necesitan OCR}; si es cero, imprime \texttt{nada que
        hacer}. Si no es cero, empieza a procesarlos directamente en el USB.
\end{itemize}

En la prueba de 50 PDFs del 4 de noviembre, los 50 quedaron con texto: el menor con 2\,902
caracteres, la mediana con 21\,168 y el mayor con 94\,607. El texto es legible, pero en
los encabezados el OCR deja frases desordenadas o con errores (por ejemplo, un membrete
partido en columnas). Conviene revisar a mano unos cuantos markdown antes de usarlos para
entrenar.

%--------------------------------------------------------------------
\section{Después del OCR}

\begin{enumerate}
  \item El markdown de 2014 ya lo genera \path{parte-sola.cmd} al final; revisar
        \path{necesitan-ocr-2014.csv} como se indica arriba.
  \item Subir a Drive \texttt{pdfs/*/2014/} y \texttt{markdown/*/2014/}, y luego marcar el
        año como subido con \texttt{lote.py --subido 2014}, según el flujo descrito en
        \texttt{testeo/README.md}. Esa orden borra los PDFs y conserva el markdown.
  \item Recordatorio: \texttt{estado.json} aún no marca 2014 como subido.
\end{enumerate}

%--------------------------------------------------------------------
\section{Cambios en el código}

Se modificó \texttt{testeo/ocr.py}:

\begin{itemize}
  \item \texttt{--output-type pdf}: no hace falta Ghostscript, que no se pudo instalar sin
        administrador. El PDF/A no se necesita para este corpus.
  \item Llama a \texttt{python -m ocrmypdf} en vez del ejecutable \texttt{ocrmypdf}, que
        solo existe con una instalación normal de pip. Así funciona con el Python portátil.
  \item Nuevo \texttt{--workers}: varios documentos en paralelo. Con el valor por defecto,
        1, se comporta como antes.
  \item La revisión previa de qué PDFs tienen texto se hace con hilos.
  \item Se corrigió un error que ya existía: combinar \texttt{--especialidad} con
        \texttt{--anio} no encontraba ningún PDF, porque el patrón buscaba un nivel de
        carpetas de más.
\end{itemize}

%--------------------------------------------------------------------
\section{Si el laboratorio no deja ejecutar los \texttt{.cmd}}

Los \texttt{.cmd} son solo atajos: llaman a PowerShell y a Python. Hay tres niveles de
respaldo, para probar en orden desde la misma consola en \texttt{X:\textbackslash
ocr-offline}. Se pasa al siguiente nivel solo cuando el anterior da un error de bloqueo
(por ejemplo, \emph{la ejecución de scripts está deshabilitada} o \emph{el administrador
ha bloqueado esta aplicación}), no cuando falla el OCR en sí.

\begin{center}
\begin{tabular}{lp{4.6cm}p{5.6cm}}
\toprule
Nivel & Cuándo & Cómo se lanza \\
\midrule
0 & Funciona todo & \texttt{parte-sola.cmd} \\
1 & El \texttt{.cmd} no, PowerShell sí & \texttt{powershell \ldots\ -File .ps1} \\
2 & PowerShell tampoco & \texttt{python.exe ocr-lab.py} \\
3 & Nada desde el USB & Copiar el kit al disco local y repetir 1 o 2 \\
\bottomrule
\end{tabular}
\end{center}

\paragraph{Nivel 1: PowerShell funciona, el \texttt{.cmd} no.}
\begin{lstlisting}
powershell -NoProfile -ExecutionPolicy Bypass -File .\ocr-lab.ps1 -Limite 50 -Especialidad civil
powershell -NoProfile -ExecutionPolicy Bypass -File .\ocr-parte.ps1 -Limpiar -Especialidad penal
powershell -NoProfile -ExecutionPolicy Bypass -File .\ocr-parte.ps1 -Limpiar -Especialidad contencioso-adm-laboral,contencioso-adm-previsional
powershell -NoProfile -ExecutionPolicy Bypass -File .\ocr-parte.ps1 -Limpiar -Especialidad revision-de-procedimiento-coactivo,contencioso-administrativo,civil,laboral,constitucional,familia-civil,comercial,familia-tutelar,familia-penal
powershell -NoProfile -ExecutionPolicy Bypass -File .\ocr-lab.ps1 -Convertir
\end{lstlisting}
Las tres líneas con \texttt{ocr-parte.ps1} son las tres tandas de \texttt{parte-sola.cmd},
y la última genera el markdown. Aquí la conversión \emph{no} se lanza sola: hay que ejecutar
esa última línea cuando terminen las tres tandas.

\paragraph{Nivel 2: PowerShell también está bloqueado.} Se usa el Python portátil
directamente, con el lanzador \texttt{ocr-lab.py}, que hace lo mismo sin \texttt{.cmd} ni
\texttt{.ps1}:
\begin{lstlisting}
portable\python\python.exe ocr-lab.py --simular --todo
portable\python\python.exe ocr-lab.py --limite 50 --especialidad civil
portable\python\python.exe ocr-lab.py --todo
\end{lstlisting}
\begin{sloppypar}
\texttt{--todo} hace las tres tandas y, al terminar, \textbf{convierte solo a markdown}, igual
que \path{parte-sola.cmd}. Con \texttt{--sin-convertir} se detiene tras el OCR, y la
conversión se lanza aparte con \texttt{ocr-lab.py --convertir}. \texttt{--simular} solo
muestra los pasos sin ejecutarlos. Admite también \texttt{--local D:\textbackslash ocr},
\texttt{--workers} y \texttt{--jobs}.
\end{sloppypar}

\paragraph{Nivel 3: no deja ejecutar nada desde el USB.} Copiar \texttt{X:\textbackslash
ocr-offline} a \texttt{C:\textbackslash ocr-offline} y repetir los niveles 1 o 2 desde
allí; los datos se siguen leyendo de \texttt{X:}. Si tampoco funciona, la PC bloquea
programas en carpetas de usuario por política del laboratorio. Eso no se resuelve desde
aquí: hay que pedir al encargado que permita ejecutar Python y Tesseract desde el USB, o
usar otra PC. No conviene intentar saltarse la política.

\paragraph{El resultado es el mismo en todos los niveles.}\sloppy Los PDFs con texto quedan en
\path|pdfs\<especialidad>\2014\|, el markdown en \path|markdown\<especialidad>\2014\|,
y los fallos en \path|pdfs\fallos-ocr.csv| y \path|markdown\necesitan-ocr-2014.csv|.
El log queda en \path{ocr-lab-2014-NOMBREPC.log}. En todos los niveles, si se corta, se
repite el mismo comando y se reanuda.

\paragraph{Estado de la prueba.} \texttt{ocr-lab.py} compila y su modo \texttt{--simular}
muestra los pasos correctos, incluida la conversión final. El flujo de \texttt{parte-sola.cmd}
se comprobó con PowerShell sustituido por un \texttt{echo}: ejecuta las tres tandas y luego
la conversión, y se detiene sin convertir si una tanda falla. Ninguno de los dos
lanzadores se ha ejecutado con OCR real.

%--------------------------------------------------------------------
\section{Qué no se ha probado}

\begin{itemize}
  \item El kit solo se ha ejecutado en la PC de casa. En el laboratorio puede haber
        bloqueos de seguridad que impidan ejecutar programas desde el USB. En ese caso,
        copiar \texttt{X:\textbackslash ocr-offline} a \texttt{C:\textbackslash
        ocr-offline} y lanzar desde allí; los datos se siguen leyendo de \texttt{X:}.
  \item El paralelismo con muchos núcleos y el ciclo completo de copiar, hacer OCR,
        devolver y borrar la copia local no se han ejecutado de principio a fin. Por eso
        la prueba de 50 va primero y la tanda A se vigila al inicio.
  \item Se asume Windows de 64 bits; los paquetes de Python incluidos son para Python 3.13.
\end{itemize}

\end{document}
